% Section 8: PDLP on CPU and GPU.
\section{PDLP: a first-order LP method on CPU and GPU}\label{sec:pdlp}

The simplex and interior-point methods need sparse factorisations, and those are hard to run
efficiently on a GPU. PDLP~\cite{applegate2021} needs only products with $A$ and $A\T$ and vector
operations, which GPUs do well. \file{lp/pdlp.py} implements the restarted, preconditioned
primal--dual hybrid gradient method (PDHG~\cite{chambolle2011}) with the enhancements of PDLP and of
its GPU descendants cuPDLP.jl~\cite{lu2023cupdlp} and cuPDLP-C~\cite{lu2023cupdlpc}. The method is
written once over an array module (NumPy or CuPy); the inner loop has a Numba backend for the CPU
and a CUDA backend for the GPU.

\subsection{The PDHG step}
For the saddle-point form \eqref{eq:saddle}, one step with primal step $\tau=\eta/\omega$ and dual
step $\sigma=\eta\omega$ (step size $\eta$, primal weight $\omega$) is
\begin{align}
  x' &= \proj_X\bigl(x-\tau(c-A\T y)\bigr),\label{eq:pdhg-x}\\
  y' &= \arg\max_{y\in Y}\; -y\T A(2x'-x) + p(y) - \tfrac{1}{2\sigma}\|y-y_k\|^2 .\label{eq:pdhg-y}
\end{align}
The dual step has a closed form. With $t=y-\sigma A(2x'-x)$, Moreau's identity gives
\begin{equation}\label{eq:moreau}
  y' = t + \sigma\,\proj_{[r^l,r^u]}(-t/\sigma) = t - \clip\bigl(t,\,-\sigma r^u,\,-\sigma r^l\bigr),
\end{equation}
so an equality row gets $y-\sigma(Ax'_{\text{extr}}-b)$, a $\ge$ row gets $\max(0,t+\sigma r^l)$, and a
free row gets~$0$. No slack columns are introduced: ranged rows, bounds and free rows are all
handled by the two projections.

\paragraph{Adaptive steps}
PDLP's step-size rule (its Algorithm~3) accepts a step when $\eta\le\bar\eta$, where
\begin{equation}\label{eq:pdlp-step}
  \bar\eta = \frac{\omega\|x'-x\|^2+\|y'-y\|^2/\omega}{2\,\bigl|(x'-x)\T A\T(y'-y)\bigr|},\qquad
  \eta_{\text{next}} = \min\Bigl(\bigl(1-(k+1)^{-0.3}\bigr)\bar\eta,\ \bigl(1+(k+1)^{-0.6}\bigr)\eta\Bigr).
\end{equation}
A rejected step is not committed and is retried with $\eta_{\text{next}}$. The first step is
$\eta_0=1/\|A\|_{\max}$. Accepted iterates enter a running average weighted by their step size.

\paragraph{Restarts and the primal weight}
Every 64 attempts the method evaluates the primal-weighted KKT error
$\sqrt{\omega\,\|r_p\|^2+\|r_d\|^2/\omega+\text{gap}^2}$ at the current iterate and at the
average, and restarts from the better of the two, following cuPDLP's criteria~\cite{lu2023cupdlp}.
It restarts when that error has fallen below $0.2$ times its value at the last restart; or below $0.8$
times that value while rising since the previous check; or when the iterations since the last restart
exceed $0.36$ of all iterations. Restarts are what give PDHG linear convergence on LPs~\cite{applegate2023}.
At a restart the primal weight is smoothed towards the ratio of dual to primal movement,
\begin{equation}\label{eq:omega}
  \omega \leftarrow \exp\bigl(\theta\log(\|\Delta y\|/\|\Delta x\|)+(1-\theta)\log\omega\bigr),\qquad \theta=0.5 .
\end{equation}

\paragraph{Preconditioning}
Ten Ruiz iterations~\cite{ruiz2001} divide each row and column by the square root of its largest
entry. One Pock--Chambolle pass~\cite{pock2011} ($\alpha=1$) then divides by the square roots of the
row and column $\ell_1$ norms. Finally the bounds and the objective are rescaled by
$\beta_b=1+\|\text{bounds}\|_2$ and $\beta_c=1+\|c\|_2$, as in cuPDLP-C~\cite{lu2023cupdlpc}, so that
$\omega$ starts near~1. The initial weight is $\omega_0=\|c\|/\|b\|$ on the scaled problem.

\paragraph{Termination and infeasibility}
The tests are PDLP's relative ones, evaluated on the \emph{unscaled} problem:
$\|r_p\|_2\le\epsilon(1+\|b\|_2)$, $\|r_d\|_2\le\epsilon(1+\|c\|_2)$ and
$|c\T x-\phi|\le\epsilon(1+|c\T x|+|\phi|)$, where $r_d$ is the part of $c-A\T y$ that the variable
bounds cannot absorb. The default is $\epsilon=10^{-4}$. Every fourth check (256 attempts), the
difference between the current iterate and the one at the previous check is projected onto the
recession cones of the bounds and tested as a Farkas certificate~\cite{applegate2024}. A dual ray
with positive objective and relative infeasibility below $10^{-8}$ proves primal infeasibility; a
primal ray with $c\T\Delta x<0$ proves unboundedness.

The 2-norm primal test alone is not safe. On \code{s250r10} a few right-hand sides are huge
($\|b\|_2\approx2\times10^5$), so the test accepted a point whose worst row was violated by about 20 in
absolute terms: the reported objective was 0.22 away from the optimum, and every bound is finite, so
the dual residual gave no warning. NIRNAY therefore accepts a candidate only if, in addition, every
row satisfies $\mathrm{viol}_i\le10\epsilon\,(1+|b_i|)$ in the original model. A candidate that
fails the row check is counted (\code{row\_check\_rejects}) and the iteration continues. On
\code{s250r10} the GPU run now ends within $3\times10^{-5}$ of the HiGHS objective, at the cost of a
longer run (\cref{tab:large}).

\subsection{Two backends}
\paragraph{CPU}
One Numba call runs a whole block of 64 attempts, in three passes over the data per attempt:
(1)~the primal step \eqref{eq:pdhg-x}; (2)~a row pass that computes $(Ax')_i$ and immediately takes
the dual step \eqref{eq:moreau} for row~$i$; and (3)~a column pass that computes $(A\T y')_j$ together
with the two reductions the step rule \eqref{eq:pdlp-step} needs. The average is updated lazily inside
the next attempt's first two passes. The passes are thread-parallel when $\nnz(A)\ge30\,000$. Below
that, on the development laptop, the fork/join cost of the thread pool outweighed the gain.

\paragraph{GPU}
Each attempt is four CUDA kernels (\cref{fig:gpu}). A \emph{commit} kernel adds the previous accepted
step to the average and takes the primal step. A \emph{rows} kernel computes $Ax'$ with a group of $G$
lanes per row, reducing with warp shuffles, and fuses the dual step. A \emph{columns} kernel computes
$A\T y'$ and fuses $\|\Delta x\|^2$ and $\Delta x\T A\T\Delta y$. A one-thread \emph{step-size}
kernel applies \eqref{eq:pdlp-step}. $G$ is the largest power of two not above the mean row (or
column) length, capped at~32. The accept/reject decision stays on the device: a rejected step is simply
not committed by the next commit kernel. A block of 64 attempts therefore needs no host round trip,
and it is captured once as a CUDA graph and replayed. The host reads one small array of scalars per
check.

\begin{figure}[tbp]
\centering
\begin{tikzpicture}[x=1cm,y=1cm]
  \node[lbl,anchor=east] at (-0.2,0) {GPU stream};
  \foreach \a/\x in {1/0,2/3.35} {
    \node[blk,fill=navy,text=white,minimum width=0.72cm,minimum height=0.7cm,font=\sffamily\tiny] (c\a) at (\x+0.36,0) {commit\\x-step};
    \node[blk,fill=sky,draw=sky,text=white,minimum width=0.72cm,minimum height=0.7cm,font=\sffamily\tiny] (r\a) at (\x+1.14,0) {rows\\$Ax'$,\,$y'$};
    \node[blk,fill=igreen,draw=igreen,text=white,minimum width=0.72cm,minimum height=0.7cm,font=\sffamily\tiny] (k\a) at (\x+1.92,0) {cols\\$A\T y'$};
    \node[blk,fill=saffron,draw=saffron,text=white,minimum width=0.72cm,minimum height=0.7cm,font=\sffamily\tiny] (s\a) at (\x+2.7,0) {step\\$\bar\eta$};
  }
  \node[font=\sffamily\small] at (7.35,0) {$\cdots$};
  \node[blk,fill=navy,text=white,minimum width=0.72cm,minimum height=0.7cm,font=\sffamily\tiny] (fl) at (8.2,0) {flush};
  \draw[decorate,decoration={brace,amplitude=5pt,mirror},muted] (0,-0.5) -- node[below=5pt,lbl]{attempt 1} (3.06,-0.5);
  \draw[decorate,decoration={brace,amplitude=5pt,mirror},muted] (3.35,-0.5) -- node[below=5pt,lbl]{attempt 2} (6.41,-0.5);
  \draw[decorate,decoration={brace,amplitude=5pt},navy] (0,0.5) -- node[above=5pt,lbl,text=navy]{one CUDA graph: 64 attempts, captured once, replayed; accept/reject decided on the device} (8.56,0.5);
  \node[blkm,minimum width=3.6cm,align=center] (host) at (11.3,0) {host check every 64\\[-1pt]{\scriptsize KKT (unscaled), restart,}\\[-1pt]{\scriptsize primal weight, Farkas rays}};
  \draw[flow] (fl.east) -- (host.west);
  \draw[flow,muted] (host.south) -- ++(0,-0.9) -| node[lbl,below,pos=0.25]{next block (same graph)} (c1.south);
\end{tikzpicture}
\caption[PDLP on the GPU]{PDLP on the GPU. Each attempt is four fused kernels; the step-size
decision never leaves the device. A block of 64 attempts is one replay of a captured CUDA graph, and
the host is involved only at the check between blocks.}
\label{fig:gpu}
\end{figure}

\subsection{Results}\label{sec:pdlp-results}
\paragraph{Netlib} PDLP at its default tolerance $10^{-4}$ reports \code{optimal} on \nPdlp{} of the
\nPdlpTot{} Netlib problems (\tlPdlp\,s limit). \nPdlpTimeLimit{} hit the time limit
(\pdlpTimeLimitNames) and \nPdlpIterLimit{} the iteration limit (\pdlpIterLimitNames). These are
the badly scaled problems that every PDLP implementation in our comparator study also fails
(\cref{sec:competitors}). A relative KKT tolerance of $10^{-4}$ does not bound the objective error by
$10^{-4}$: \nPdlpTight{} of the problems are within $10^{-4}$ of the HiGHS objective,
\nPdlpLoose{} within $10^{-3}$, and only \nPdlpStrict{} within the $10^{-6}$ used to score the other
engines. The worst objective errors among the ``optimal'' runs are \pdlpWorst. The median run takes
\medItPdlp{} iterations (the longest \maxItPdlp). The shifted geometric mean time is \sgmPdlp\,s
against \sgmPdlpRef\,s for HiGHS. Netlib problems are small and hard, and PDLP is the wrong tool for
them. It is there to show correctness, not speed.

\ifnum\haveLarge=1
\paragraph{Large LPs, CPU against GPU}
The regime PDLP is built for is large, structured LPs. \Cref{tab:large,fig:large} show the
\nLarge{} large instances run so far\ifnum\nLarge<\nLargeList{} (the sweep over \nLargeList{} files is
still running)\fi, from the Mittelmann benchmark sets~\cite{mittelmann}, all at relative tolerance
\largeTol{}, on an \largeGpuHw{} against the CPU backend on the laptop's six cores (12 threads). On
the \nLargeBothOk{} instances both backends finish, the GPU-over-CPU speed-up ranges from
\largeSpeedMin{} to \largeSpeedMax{} (geometric mean \largeSpeedGm; a value below~1 means the GPU was
slower). On \nLargeGpuOnly{} more
(\largeGpuOnlyNames) only the GPU finishes within the \largeLimit\,s limit, and \nLargeNeither{}
finished on neither. HiGHS in the same table runs its interior-point method without crossover at its
default $10^{-8}$ tolerances, which gives the reference objective. That is a far more accurate answer,
so the comparison with the PDLP times is not like for like. HiGHS's interior point at the same
$10^{-4}$ tolerance (last column) is the closer comparison. Where both finish, the ratio of the
high-accuracy HiGHS time to the GPU time ranges from \largeVsHighsMin{} to \largeVsHighsMax.

\begin{table}[ht]
\centering\small
\caption[PDLP on large LPs: CPU against GPU]{PDLP on large LPs (LP relaxations for the MIPLIB
files): wall time of \nirnay{} PDLP on the CPU (Numba, 12 threads) and on the GPU (CuPy), the
relative objective error of the GPU run against the HiGHS reference, and HiGHS's interior point
at its default tolerance (reference) and at the PDLP tolerance. $^\dagger$: stopped at the time
limit; ``$>$'' is then a lower bound on the speed-up. $^\ddagger$: HiGHS did not report
\code{Optimal}.}
\label{tab:large}
\setlength{\tabcolsep}{4pt}
\begin{tabular}{@{}l rrr rrr r rr@{}}
\toprule
 & & & & \multicolumn{3}{c}{\nirnay{} PDLP} & & \multicolumn{2}{c}{HiGHS IPM (s)}\\
\cmidrule(lr){5-7}\cmidrule(l){9-10}
Instance & $m$ & $n$ & nnz & CPU (s) & GPU (s) & $\times$ & rel.\ err & $10^{-8}$ & $10^{-4}$\\
\midrule
\rows{data/large_pdlp.tex}
\bottomrule
\end{tabular}
\end{table}

\begin{figure}[ht]
\centering
\begin{tikzpicture}
\begin{axis}[
  width=0.95\linewidth, height=6cm,
  ybar, bar width=5pt, ymode=log, log origin=infty,
  ylabel={wall time (s)},
  xtick=data, xticklabels from table={data/large_pdlp.dat}{name},
  xticklabel style={font=\ttfamily\scriptsize,rotate=30,anchor=north east},
  enlarge x limits=0.06,
  ymajorgrids, legend style={at={(0.5,1.03)},anchor=south}, legend columns=3,
  ymin=0.1, ymax=1000,
]
\addplot[fill=muted!60,draw=muted] table[x=x,y=cpu] {data/large_pdlp.dat};
\addplot[fill=saffron,draw=saffron] table[x=x,y=gpu] {data/large_pdlp.dat};
\addplot[fill=navy,draw=navy] table[x=x,y=highs] {data/large_pdlp.dat};
\legend{{\nirnay{} PDLP (CPU)},{\nirnay{} PDLP (GPU)},{HiGHS IPM $10^{-8}$}}
\end{axis}
\end{tikzpicture}
\caption[PDLP on large LPs]{Wall time on the large LPs (log scale), from
\file{results/large\_pdlp\_gpu.csv}. The PDLP runs stop at relative tolerance $10^{-4}$;
the HiGHS reference solves to $10^{-8}$.}
\label{fig:large}
\end{figure}
\fi
